\scanpage{114}
% archive foliation: 113

\deleted{B) Cas général. On se réduit au}

Si $F$ est quelconque, on écrit que $F$ est un sous-\deleted{objet \inserted{\deleted{cst tordu}} d'un quotient \inserted{\deleted{cst tordu}} d'un $R^d f_*(\mathbb Q_{lY})$}, où $f:Y\longrightarrow X$ morphisme projectif. [Utilisant Lefschetz \deleted{et 2.1 d) \unsure{$f,g'$} \unsure{$j\geq$ dim relative}} \deleted{faible, on peut supposer que les fibres de $f$} \inserted{fort} si $j<d=\dim$ relative de $Y/X$, et remplaçant \deleted{sont de dim $j$}. $Y$ par $Y\times_X\mathbb P_X^{j-d}$ si $j>d$, on peut supposer $j=d$]

\margin{\deleted{\unsure{et} Lefschetz fort}}
\margin{voir si cette réduction est vraiment utile \ldots}

La suite spectrale
\[
R^\alpha_!(fg)(\mathbb Q_l)\Longleftarrow E^{pq}_2=R^p_!f((R^q_!g(\mathbb Q_{lX})))
\]
dégénère, \uline{par Deligne}, donc
\[
R^i_!f(F)=E^{ij}_2
\]
est $\simeq$ quotient \inserted{cst tordu} d'un ss faisceau ct tordu d'un $R^{i+j}_!\unsure{g|fg}(\mathbb Q_{lY})$.

\margin{\uuline{utilise Lefschetz fort}}
\margin{[N.B. Voir si Deligne ne peut se remplacer par \uline{\unsure{factorisation} de Artin} sur $X\longrightarrow S$, \unsure{puis} ramener au cas de $n=1$\ldots]}
\margin{semi-simplicité utilisée!}

La conclusion en résulte, en utilisant le cas précédent pour $Y\longrightarrow S$, de dim relative $\leq n+j$ \ldots

\uline{Corollaire} \unsure{5.2|5.1}. Si la partie ``\uuline{semi-simplicité}'' des conjectures de Tate est vérifiée, la conclusion de 5.1. est vraie dès que l'on suppose seulement que $F$ est eff$^{t}$ admissible de poids $j$. Par suite, si on suppose $F$ eff$^{t}$ admissible de poids $\leq\deleted{\unsure{k}}\inserted{j}$, on a un énoncé correspondant, \deleted{De} $R^i\unsure{f_*|f_!}(F)$ est \inserted{eff. admissible,} de poids $\leq\deleted{\unsure{\operatorname{Inf}}}i+j$, \deleted{\unsure{de niveau} $\unsure{i+j}$} \inserted{niveaux $\leq\operatorname{Inf}(i+j,2n-i+j)$}. Si $i\geq n$, on trouve des poids entre $2(i-n)$ et $(i+j)$, correspondant aux niveaux $0,1,\ldots,2n-i+j$.

On se ramène en effet \inserted{par dévissage et déf;} au cas où $F$ est \uline{spécial} de poids $j$, et par th. de chgt. de base propre au cas où $S$ de type fini sur $\unsure{\operatorname{Sp}|\operatorname{Spec}}\mathbb Z$, et

\margin{alors par semi-simplicité $F$ est \unsure{stt} spécial}

\uline{Cor.} 5.3 (sans conj. type Tate) \deleted{\unsure{on a}} Si $F$ est eff. admissible, Alors $R^i_!f(F)$ est eff$^{t}$ admissible.
